\section{引言}

接吻数是离散几何中历史最悠久的问题之一。1694年，牛顿与格雷戈里讨论过：
三维空间中，究竟能有12个还是13个等大的球同时接触一个中心球？
答案12直到两个多世纪后才得到证明\cite{musin}。
这个问题看似简单，实际却极为困难。三百多年过去，
精确接吻数仍只在1、2、3、4、8和24维得到确定。
其中，8维与24维的最优构型分别来自$E_8$格和Leech格，
接触数为240和196560~\cite{conway1999sphere,musin}。
在大多数其他维度，已知构造与上界之间仍存在显著差距。

下界的进步长期依赖明确的数学结构。格与纠错码提供了大型向量族，
使大量点对的关系能够同时受到控制\cite{conway1999sphere,macwilliams}。
由此发展出的层叠格、格截面与扩张、编码球面构型，以及将高维母构型
与新增坐标组合的提升方法，可产生数十万乃至数千万个点
\cite{laminated,antipode,ers,cjkt,kallal}。
但结构的优势也带来了困难：当一个构型受到严密约束时，改进它可能需要
重新组织整组点，而非简单插入一个新点。搜索空间随组合选择迅速增长，
真正允许改进的几何自由度，却未必能从原有表示中直接看出。

人工智能正在改变探索这些空间的方式。通用编码智能体AlphaEvolve在11维
发现了593点接吻构型，将此前的592点下界提高了一点\cite{alphaevolve}。
PackingStar将接吻数问题表述为双参与者的矩阵补全博弈，
以博弈论强化学习搜索Gram矩阵空间，在25至31维报告新构造，
并发现数千种候选几何结构\cite{packingstar-discovery}。
EinsteinArena允许自主智能体提交解、查看中间结果并继承其他智能体的发现；
这一协作过程将11维下界从593提高到604~\cite{einsteinarena}。
随后，Station在开放式多智能体环境中给出了精确的604点构型，
并发展出结构性定理\cite{station}。
这些工作展示了人工智能在几何搜索与协作式数学发现中的能力。

更深层的问题也由此显现。找到一个更大的构型，只是数学发现的一部分：
它为什么存在？什么结构自由度使改进成为可能？这种机制能否推广？
怎样将结果落实为证明或可独立核验的证书？
计算搜索通常从给定的表示出发，例如坐标、Gram矩阵、编码或固定的优化目标。
开放式数学研究还需要在原有表示成为限制时改变表示，
提出新的中间问题，连接看似无关的结构，并发展能够证明构造成立的论证。

自主科研系统为这类发现提供了一条途径。基于大语言模型的智能体可以
持续推理、计算并使用工具，根据中间结果调整研究策略。
Qiushi Engine此前已在真实光学平台上开展端到端自主科学研究，
将理论推理、计算与实验结合起来，完成长程探索\cite{qiushi-optics}。
本研究考察这一范式能否进入开放式数学研究：
自主系统是否能够超越固定数学表述下的优化，发现新的表示、构造与证明？

我们报告Qiushi Engine对接吻数问题的自主研究，得到十九个维数的新下界：
\[
25,\quad27,\quad32\text{—}39,\quad43,\quad45,\quad49\text{—}55.
\]
代表性结果包括
\[
\begin{aligned}
K(25)&\ge197580,&K(27)&\ge201567,&K(38)&\ge591900,\\
K(43)&\ge2553792,&K(45)&\ge7380090,&K(55)&\ge53301140.
\end{aligned}
\]
这些改进并非同一优化程序的重复运行，而来自一系列不同的结构发现：
整个接触层的协同运动、分离方向删除与复用的带标签提升、
利用共享删除代价的联合支持交换、通过Hadamard转移得到的符号码替换、
从另一格壳选取相容方向、扩充格直线类并组合低重叠等距像，
以及无需枚举整个母壳便能确定或控制格截面点数的球面设计矩恒等式与不等式。

这些构造具有共同的认识：在一种表示中看似饱和的构型，
换一种表示后仍可能具有可利用的自由度。
25维让552个接触层点整体移动，保持其内部Gram矩阵，为一个新点腾出空间。
27维的带标签提升将方向复用与删除分开；最终两层构造保持第一层不变，
联合选择第二层方向与标签，从而扩大点集。
32至37维及39维通过联合交换与单独选择符号，
利用逐个插入评分和完整符号束所遮蔽的改进空间。
38维将搜索从Leech极小壳转向第二壳层，增加288个点。
43维和45维把数百万个极小向量的计数转化为有限投影签名上的重数问题，
由精确球面设计矩施加约束。
48维中同一个扩充的直线类，再与低重叠等距像组合，
则产生49至55维的全部七项构造。

Qiushi Engine自主完成了构造搜索、数学分析与计算验证。
每项发现均表述为明确的几何或组合条件，各项下界由有限见证支持，
通过精确整数与有理数运算、格与码的恒等式、严格区间算术及有限组合检查建立。
其中，25维使用向外舍入的区间算术认证剩余的严格不等式
\cite{johansson-arb}。
这些数学论证与可复现证书使构造能够被独立检查。

这些成果表明，人工智能对数学研究的作用可以超越数值搜索的加速。
面对庞大而有结构的搜索空间，关键任务可能是找到使进展成为可能的数学表示。
能够持续进行探索、结构分析与严格验证的自主科研系统，
为大规模开展这类发现提供了新的途径。

\subsection{球面表述}

接吻数$K(n)$是欧氏$n$维空间中能够同时接触中心单位球、
而彼此内部不相交的单位球的最大数量。
将外围球心除以二，便得到等价的球面问题：
\[
K(n)=\max\{|X|:X\subset S^{n-1},\
\langle x,y\rangle\le\tfrac12\quad(x\ne y)\}.
\]
不同方向之间的夹角至少为$60^\circ$。

\subsection{主要结果}

\begin{theorem}\label{thm:main}
对表\ref{tab:results}中的每个维数 $n$，存在含 $N_n$ 个点的接吻构型，其中 $N_n$ 为第二列的下界。因此 $K(n)\ge N_n$。
\end{theorem}

本报告中的构造给出下列下界。比较列采用2026年9月24日至27日核对的公开固定来源。

\begin{longtable}{rrrrl}
\caption{显式构造下界及其公开比较值。}\label{tab:results}\\
\toprule
维数 & 下界 & 公开比较值 & 增加 & 来源\\\midrule
\endfirsthead
\toprule
维数 & 下界 & 公开比较值 & 增加 & 来源\\\midrule
\endhead
\bottomrule\endfoot
25 & 197580 & 197579 & 1 & \cite{kissingnumbers} \\
27 & 201567 & 201566 & 1 & \cite{lindow27} \\
32 & 347584 & 346944 & 640 & \cite{brouwer} \\
33 & 363968 & 362048 & 1920 & \cite{brouwer} \\
34 & 384196 & 381124 & 3072 & \cite{brouwer} \\
35 & 409676 & 409548 & 128 & \cite{brouwer} \\
36 & 484760 & 484568 & 192 & \cite{brouwer} \\
37 & 498024 & 496232 & 1792 & \cite{brouwer} \\
38 & 591900 & 591612 & 288 & \cite{kissingnumbers} \\
39 & 763668 & 756116 & 7552 & \cite{kissingnumbers} \\
43 & 2553792 & 2545056 & 8736 & \cite{sunwang-sections,crosssections} \\
45 & 7380090 & 7379838 & 252 & \cite{sunwang-sections,crosssections} \\
49 & 52430156 & 52430140 & 16 & \cite{kissingnumbers} \\
50 & 52458468 & 52458418 & 50 & \cite{kissingnumbers} \\
51 & 52500930 & 52500816 & 114 & \cite{kissingnumbers} \\
52 & 52585772 & 52585516 & 256 & \cite{kissingnumbers} \\
53 & 52698696 & 52698222 & 474 & \cite{kissingnumbers} \\
54 & 52923774 & 52922906 & 868 & \cite{kissingnumbers} \\
55 & 53301140 & 53299730 & 1410 & \cite{kissingnumbers} \\
\end{longtable}

32维比较值采用Brouwer码表所载Lysenst\o en的1671字码\cite{brouwer}，
代入ERS公式得到$2^{17}+128\cdot1671+2\cdot32\cdot31=346944$。
研究实际使用的搜索种子仍是Echols的1667字码\cite{echols}。


\subsection{构造原理}

十九项下界来自七类相互联系的机制。每一类都改变具体的数学对象，并有对应的
相容性或计数论证。

\paragraph{25维的协同移动。}
552个接触层点具有固定的横向坐标和共同的纬度。让纬度与新增坐标共同变化，
可以保持全部内部内积；水平投影始终位于旧壳点之间的弦上，因此凸性统一控制
未动赤道。两个提升点的修复使新点得以加入，将197579增至197580。
这里利用的是保留点集的几何运动，而非在固定候选池中多选一个点。

\paragraph{27维的方向复用与标签联合选择。}
一个方向提升后可贡献多个副本，却只需删除一个赤道点；尾几何决定允许的复用。
四三角形构造达到指定十二标签模型的容量。更大的两层构造保留第一层，
联合重选第二层方向及其标签，以311个方向取代310个方向，得到201567点。
模型容量与最终两层点数分别刻画同一提升几何的不同部分。

\paragraph{32—34、37、39维的共享删除代价。}
对于一组彼此相容的新支持 $Y$，必须删除的旧支持构成并集 $F(Y)$，
净收益为 $|Y|-|F(Y)|$。即使每次单独插入都不能增点，多个候选仍可能通过
共享删除而共同成功。对应的新支持族大小为1676、1803、1960、2845和3383。
支持交叠条件与配套稠密码的距离条件共同建立完整的分层构型。

\paragraph{35—37维的符号设计转移。}
选取与外围构型相容的坐标区域，可以在其中构造更大的符号替换。
对坐标配对后施加Hadamard变换，将横截权四区组转为权八符号向量，
范数与内积同比缩放。所得局部码分别有256、576和512个点；
第一项达到十一坐标模型的容量，第二项可由两组一因子分解中的四种颜色直接构造。
37维则把十三个额外支持带来的1664点与符号替换的128点结合。

\paragraph{38维的跨壳赤道补充。}
基底提升使用了全部Leech极小方向，但仍留下空赤道。在同一整数Leech格中，
平方范数分别为48和32的向量 $v,u$ 满足 $\|v\pm u\|^2\ge32$，
因而 $|v\cdot u|\le24$。这一关系将全部旧壳点从子集选择的约束中消去，
只需寻找彼此相容的新直线。144条直线贡献288个点，得到591900点构型。

\paragraph{49—55维的像族并集。}
$P_{48p}$ 中扩充到7077条极小直线的相容类，经等距变换产生更多相容类。
决定增益的是像的并集：重复直线只分配一次。若并集包含 $S$ 条直线，
配以 $t$ 条不同的反极尾直线，构型点数为 $52416000+2S+2t$。
同一提升原理给出七项下界；各维度的像族分别指定，并将母类扩充和像族选择
各自带来的增益明确分开。

\paragraph{43、45维的目标矩计数。}
48维极值格的极小壳构成球面11设计，其矩约束所有可容许投影纤维的重数。
对任意含$\sqrt3E_8$的48维极值偶幺模格，有理恒等式精确确定
与指定$D_5$子截面及其各Weyl共轭子截面正交的极小向量数均为2553792。
45维的矩方程与范数八截面约束共同给出精确恒等式
\[
f(e_1)+f(e_2)+f(e_3)=7377408+9f(-2,-2,3).
\]
最后一个纤维对应锚图中的邻集之差，大小为 $368-c$，其中 $c$ 是所选
非邻接顶点对的共同邻点数。因此投影构型恰有 $7380720-9c$ 个点；
给定锚图在 $c=70$ 时达到最大值7380090。这两个维数都由有限的几何约束
确定大规模截面的点数。

\subsection{数学背景与直接前驱}

Cohn、Jiao、Kumar、Torquato的Leech提升框架\cite[定理7.5]{cjkt}
及后续研究，将母方向与低维尾向量组合。四三角形构造采用PackingStar提供的
496点块\cite{packingstar,packingstar-data}。
25维变形以Kravatskiy的197579点构型为起点，27维两层输入来自Lindow的
201566点构造\cite{kissingnumbers,lindow27}。38维的空赤道提升与48维的
7069条直线比较类也来自文献\cite{kissingnumbers}。
这里的新对象是协同运动及修复、第二层分配、新壳直线，以及扩充的相容类和像族。

码构造采用Edel、Rains、Sloane的分层方法\cite{ers}、
Cheng--Sloane稠密码\cite{chengsloane}及常重码构造与维护表
\cite{bsss,brouwer,echols}。局部替换的直接前驱是Takhanov--Yun的
自由区域方法\cite[第V节]{signed-johnson}；51区组的权四来源是经典的
Kalbfleisch--Stanton设计\cite{kalbfleisch-stanton}。
本研究给出的新支持、坐标配对、符号选择和外围嵌入产生更大的构型。

格截面与antipode构造源于Conway--Sloane\cite{laminated,antipode}。
43、45维的公开比较值来自Sun--Wang的计数\cite{sunwang-sections,crosssections}；
Dorofeev、Sun、Wang对嵌入 $E_8$ 的分析\cite{dorofeev}提供母截面几何。
球面设计理论\cite{dgs,nebe-designs}给出矩恒等式；
本文选择能够控制目标点数的统计量与有限见证，得到所述构型。

\subsection{阅读结构与开放构造}

报告从决定构造的几何关系出发，逐步说明具体有限对象如何形成；
附录给出坐标、检查方法与研究脉络。配套仓库包含中英文逐维叙述、原始数学数据
和程序，读者可以重建对象、核验各证明的有限条件，并进一步改变构造中的选择。
